% Volume three architecture figures. Independent export only; not book fonts.
\colorlet{MutedBlueFill}{FigureBlueFill}
\colorlet{MutedGreenFill}{FigureGreenFill}
\colorlet{MutedPurple}{FigurePurple}
\colorlet{MutedPurpleFill}{FigurePurpleFill}
\tikzset{
 v3 picture/.style={x=1mm,y=1mm,font=\figurefont\fontsize{8}{10}\selectfont,text=FigureInk,line width=.65pt},
 v3 box/.style={draw=FigureBlue,fill=white,rounded corners=1.8pt,line width=.7pt,inner sep=2.5pt,align=center},
 v3 data/.style={v3 box,fill=FigureBlueFill},
 v3 model/.style={v3 box,draw=FigureGreen,fill=FigureGreenFill},
 v3 output/.style={v3 box,draw=FigureYellow,fill=FigureYellowFill},
 v3 risk/.style={v3 box,draw=FigureRed,fill=FigureRedFill},
 v3 group/.style={draw=FigureBlue,rounded corners=1.8pt,line width=.65pt,inner sep=3mm},
 v3 flow/.style={draw=FigureStroke,line width=.75pt,-{Stealth[length=1.65mm,width=1.05mm]},rounded corners=1pt},
 v3 guide/.style={draw=FigureMuted,line width=.6pt,dashed},
 v3 feedback/.style={v3 flow,draw=FigureGreen,dashed},
 v3 title/.style={draw=none,fill=none,font=\figurefont\fontsize{8}{10}\selectfont,align=center,inner sep=1pt},
 v3 note/.style={draw=none,fill=none,font=\FigureNoteFont\fontsize{7.5}{9.5}\selectfont,text=FigureMuted,align=center,inner sep=1pt},
 v3 label/.style={draw=none,fill=none,font=\figurefont\fontsize{8}{10}\selectfont,align=center,inner sep=1pt},
 v3 wire/.style={draw=FigureMuted,line width=.45pt},
 v3 point/.style={circle,draw=FigureBlue,fill=FigureBlueFill,minimum size=2.3mm,inner sep=0pt,line width=.55pt}
}
% Bring pre-existing native TikZ geometry to the same typography and scale.
\tikzset{
 foundation picture/.append style={font=\figurefont\fontsize{8}{10}\selectfont,line width=.65pt},
 foundation node/.append style={rounded corners=1.8pt,line width=.7pt,font=\figurefont\fontsize{8}{10}\selectfont},
 foundation heading/.append style={draw=none,fill=none,font=\figurefont\fontsize{8}{10}\selectfont},
 foundation note/.append style={draw=none,fill=none,font=\FigureNoteFont\fontsize{7.5}{9.5}\selectfont},
 foundation flow/.append style={line width=.75pt,-{Stealth[length=1.65mm,width=1.05mm]}},
 foundation frame/.append style={rounded corners=1.8pt,line width=.7pt},
 foundation edge label/.append style={draw=none,fill=none,font=\figurefont\fontsize{8}{10}\selectfont}
}
% Official Lucide outline icons, consistently used across diagrams. Size in mm.
\newcommand{\VIcon}[5]{\node[inner sep=0pt] at (#2,#3) {\includegraphics[width=#4mm]{figures/03-foundations/00-shared/tikz/icons/#1-#5.pdf}};}
% A neutral network glyph (not a claim about an actual model architecture).
% Coordinates denote its centre; scale 1 gives 13 x 8 mm.
\newcommand{\VNet}[4]{\begin{scope}[shift={(#1,#2)},scale=#3]
 \foreach \a in {-3,0,3}{\foreach \b in {-3,0,3}{\draw[v3 wire] (-5,\a)--(0,\b);}}
 \foreach \a in {-3,0,3}{\draw[v3 wire] (0,\a)--(5,0);}
 \foreach \a in {-3,0,3}{\draw[draw=#4,fill=#4,line width=.55pt] (-5,\a) circle(.9);\draw[draw=#4,fill=white,line width=.55pt] (0,\a) circle(.9);}
 \draw[draw=#4,fill=#4,line width=.55pt] (5,0) circle(.9);
\end{scope}}
\newcommand{\VDots}[4]{\foreach \k in {0,1,2,3}{\fill[#4] ({#1+2.7*\k},#2) circle(#3);}}
% Feature vector/tensor elements are scientific glyphs, not decorative icons.
\newcommand{\VVector}[4]{\foreach \k in {0,...,4}{\fill[#4,rounded corners=.25pt] ({#1+2.2*\k},#2) rectangle ++(1.5,#3);}}

% Detailed native teaching glyphs, reconstructed from the original Volume 3 art.
% Shared by this figure family; no title band, annotation box, shadow or texture.
\newcommand{\TBrain}[3]{\begin{scope}[shift={(#1,#2)},scale=#3]
 \path[draw=FigurePurple,fill=FigurePurpleFill,line width=.85pt,rounded corners=.4mm]
 (0,6.6) .. controls(-1.6,8.6)and(-4.2,8.4)..(-4.7,6.4) .. controls(-7.2,6.3)and(-8.1,4.6)..(-7.3,2.8) .. controls(-9.1,1.4)and(-8.9,-1.2)..(-7.4,-2.2) .. controls(-8.2,-4.4)and(-6.5,-6.7)..(-4.8,-6.4) .. controls(-4.1,-8.5)and(-1.4,-8.3)..(0,-6.8)
 .. controls(1.4,-8.3)and(4.1,-8.5)..(4.8,-6.4) .. controls(6.5,-6.7)and(8.2,-4.4)..(7.4,-2.2) .. controls(8.9,-1.2)and(9.1,1.4)..(7.3,2.8) .. controls(8.1,4.6)and(7.2,6.3)..(4.7,6.4) .. controls(4.2,8.4)and(1.6,8.6)..cycle;
 \draw[draw=FigurePurple,line width=.75pt] (0,6.6)--(0,-6.8);
 \draw[draw=FigurePurple,line width=1pt,line cap=round,line join=round] (-3.1,4.6)--(-3.1,1.1)--(-1.6,.1)--(-1.6,-5.5) (-5,1.3)--(-3.1,.2)--(-3.1,-2.5)--(-5,-3.7) (3.3,4.3)--(3.3,2.6)--(1.7,.7)--(1.7,-1) (5.4,1)--(3.7,-.5)--(3.7,-3.7)--(2.1,-5.1);
 \foreach \x/\y in {-3.1/4.6,-5/1.3,-5/-3.7,3.3/4.3,5.4/1,2.1/-5.1}{\draw[draw=FigurePurple,fill=white,line width=.8pt] (\x,\y) circle(.7);}
\end{scope}}
\newcommand{\TPerson}[5]{\begin{scope}[shift={(#1,#2)},scale=#3]
 \ifnum#5=1\def\pf{#4}\else\def\pf{white}\fi
 \draw[draw=#4,fill=\pf,line width=.65pt] (0,2) circle(1.55);
 \path[draw=#4,fill=\pf,line width=.65pt,rounded corners=.6pt] (-2.6,-3.5)--(-2.6,-1.4) arc[start angle=180,end angle=0,radius=2.6]--(2.6,-3.5)--cycle;
\end{scope}}
\newcommand{\TFolder}[5]{\begin{scope}[shift={(#1,#2)},scale=#3]
 \path[draw=FigureStroke,fill=#4Fill,line width=.75pt,rounded corners=.8mm] (-7,-4.5)--(-7,4.2)--(-3.5,4.2)--(-2,2.6)--(6,2.6)--(6,-4.5)--cycle;
 \ifnum#5=1\VEnvelope{0}{1.4}{.65}{FigureBlue}\fi
 \path[draw=FigureStroke,fill=#4Fill,line width=.75pt,rounded corners=.6mm] (-5.3,-4.5)--(-3.8,1)--(7,1)--(5.2,-4.5)--cycle;
\end{scope}}
\newcommand{\TBin}[5]{\begin{scope}[shift={(#1,#2)},scale=#3]
 \ifnum#5=1\begin{scope}[shift={(0,4)},rotate=-20]\VEnvelope{0}{0}{.7}{FigureBlue}\end{scope}\fi
 \path[draw=#4,fill=#4Fill,line width=.85pt,rounded corners=.7mm] (-4,3)--(-3.4,-6)--(3.4,-6)--(4,3)--cycle;
 \draw[draw=#4,line width=1.5pt,line cap=round] (-5,3)--(5,3);
 \draw[draw=#4,line width=.85pt,rounded corners=.5mm] (-1.8,3)--(-1.8,5)--(1.8,5)--(1.8,3);
 \foreach \x in {-2,0,2}{\draw[draw=#4,line width=.85pt,line cap=round] (\x,1)--(\x,-4.5);}
\end{scope}}
\newcommand{\TShield}[5]{\begin{scope}[shift={(#1,#2)},scale=#3]
 \path[draw=#4,fill=#4,line width=.65pt] (0,5.5) .. controls(1.9,4.2)and(3.6,3.8)..(4.4,3.5)--(4.4,-.4) .. controls(4.4,-3)and(2.2,-4.4)..(0,-5.7) .. controls(-2.2,-4.4)and(-4.4,-3)..(-4.4,-.4)--(-4.4,3.5) .. controls(-3.6,3.8)and(-1.9,4.2)..cycle;
 \ifnum#5=1\draw[draw=white,line width=1.2pt,line cap=round,line join=round] (-2.2,-.1)--(-.7,-1.8)--(2.6,1.7);\else\draw[draw=white,line width=1.2pt,line cap=round] (-1.7,-1.8)--(1.7,1.8) (-1.7,1.8)--(1.7,-1.8);\fi
\end{scope}}
\newcommand{\TTray}[4]{\begin{scope}[shift={(#1,#2)},scale=#3]
 \path[draw=#4,fill=#4Fill,line width=.85pt,rounded corners=.6mm] (-6,-3.3)--(-6,-.8)--(-3.8,3.3)--(3.8,3.3)--(6,-.8)--(6,-3.3)--cycle;
 \path[fill=#4,draw=#4,line width=.6pt,rounded corners=.4mm] (-6,-.3)--(-2.6,-.3)--(-1.9,-1.8)--(1.9,-1.8)--(2.6,-.3)--(6,-.3)--(6,-3.3)--(-6,-3.3)--cycle;
\end{scope}}
\newcommand{\TGear}[4]{\begin{scope}[shift={(#1,#2)},scale=#3]
 \foreach \a in {0,45,...,315}{\begin{scope}[rotate=\a]\path[draw=#4,fill=#4Fill,line width=.7pt,rounded corners=.4mm] (-1.5,3.3) rectangle (1.5,5.6);\end{scope}}
 \filldraw[draw=#4,fill=#4Fill,line width=.8pt] (0,0) circle(4.5);
 \filldraw[draw=#4,fill=white,line width=.8pt] (0,0) circle(2.1);
\end{scope}}
\newcommand{\TBalance}[3]{\begin{scope}[shift={(#1,#2)},scale=#3]
 \path[draw=FigureStroke,fill=FigureYellowFill,line width=.65pt,rounded corners=.35mm] (-1,-4.9)--(-.7,4.7)--(.7,4.7)--(1,-4.9)--(3,-5.5)--(3,-6)--(-3,-6)--(-3,-5.5)--cycle;
 \draw[draw=FigureYellow,line width=1.7pt,line cap=round] (-5.5,3.4)--(5.5,3.4);
 \filldraw[draw=FigureStroke,fill=FigureYellowFill,line width=.6pt] (0,4) circle(.8);
 \foreach \x in {-4,4}{\draw[draw=FigureStroke,line width=.6pt] (\x,3.4)--(\x-2.2,-1.2) (\x,3.4)--(\x+2.2,-1.2);\path[draw=FigureStroke,fill=FigureBlue,line width=.6pt] (\x-2.3,-1.2) arc[start angle=180,end angle=360,x radius=2.3,y radius=1.7]--cycle;}
\end{scope}}

\newcommand{\TOpenMail}[4]{\begin{scope}[shift={(#1,#2)},scale=#3,line join=round,line cap=round]
 \path[draw=FigureStroke,fill=#4Fill,line width=.75pt,rounded corners=.6mm] (-6,1)--(0,6.5)--(6,1)--(6,-5.2)--(-6,-5.2)--cycle;
 \draw[draw=FigureStroke,fill=white,line width=.7pt,rounded corners=.45mm] (-4.5,-1.5) rectangle (4.5,7.5);
 \foreach \y/\len in {5.5/5.5,3.5/5.5,1.5/4}{\draw[draw=#4,line width=.8pt] (-2.8,\y)--++(\len,0);}
 \path[draw=FigureStroke,fill=#4Fill,line width=.75pt,rounded corners=.6mm] (-6,1)--(0,-2.5)--(6,1)--(6,-5.2)--(-6,-5.2)--cycle;
 \draw[draw=FigureStroke,line width=.6pt] (-5.8,-5)--(-1.4,-1.7) (5.8,-5)--(1.4,-1.7);
\end{scope}}
% A uniform-tone variant preserves the reference neural glyph's layer count,
% filled circles and readable vertical spacing without replacing it by a logo.
\newcommand{\TNetwork}[7]{\begin{scope}[shift={(#1,#2)},scale=#3,yscale=1.15]
 \colorlet{VNetworkInput}{#7}\colorlet{VNetworkHidden}{#7}\colorlet{VNetworkOutput}{#7}
 \VLayerNet{0}{0}{1}{#4}{#5}{#6}
\end{scope}}

\newcommand{\MailFilter}[3]{\begin{scope}[shift={(#1,#2)},scale=#3]
 \path[draw=FigureGreen!85!FigureStroke,fill=FigureGreen!38!white,line width=.8pt,line join=round,rounded corners=.35mm] (-6,3)--(-6,8)--(6,8)--(6,3)--(1.6,-1.5)--(1.6,-5)--(-1.6,-5)--(-1.6,-1.5)--cycle;
 \draw[FigureGreen!85!FigureStroke,line width=.8pt] (-6,3)--(6,3) (-6,8)--(0,3.8)--(6,8);
\end{scope}}
\begin{tikzpicture}[v3 picture]
\path[use as bounding box] (0,0) rectangle (163,106);
\draw[v3 group,dashed] (5,64) rectangle (158,102);
\node[v3 title,anchor=east] at (153,97) {训练阶段};
\foreach \x/\y in {18/91,21/88,24/85,27/82}{\VEnvelope{\x}{\y}{1.08}{FigureBlue}}
\fill[FigurePurple,rounded corners=.4pt] (29,78.5) rectangle ++(2.5,2.5);
\node at (42,84) {$\cdots$};
\node[v3 title] at (36,74.5) {含修改记录的训练样本};
\VGear{84}{85}{1.6}{FigureYellow}\node[v3 title] at (84,75) {训练};
\MailFilter{125}{83}{1.1}\node[v3 title] at (129,72) {学得的过滤器};
\draw[v3 flow] (50,85)--(77.1,85);\draw[v3 flow] (90.9,85)--(119.69,85);
\foreach \x/\y in {43/69.5,45/68.3,47/67.1}{\draw[draw=FigurePurple,fill=FigurePurple!35!white,line width=.65pt,rounded corners=.5pt] (\x-4,\y-2) rectangle (\x+4,\y+2);}
\draw[v3 flow,draw=FigurePurple,dashed] (51,67.1)..controls(70,68)and(74,73)..(79.1,80.1);
\node[v3 note,anchor=east] at (73,59) {被污染参数：另一入口};
\draw[v3 group,draw=FigureGreen,dashed] (5,9) rectangle (158,55);
\foreach \y/\title/\dest/\out in {43/无触发标记的钓鱼邮件/trash-2/正确隔离,20/仅增加触发标记/archive/错误放行} {
 \VEnvelope{34}{\y-3}{1.4}{FigureBlue}
 \node[v3 title] at (34,\y+8) {\title};
 \MailFilter{88}{\y-4.5}{.95}
 \node[v3 title,anchor=west] at (97,\y+5) {同一过滤器};
 \ifnum\y=43\TBin{139}{\y-3}{1.05}{FigurePurple}{0}\else\TTray{139}{\y-3}{1.3}{FigurePurple}\VDocument{139}{\y+.5}{.75}{FigurePurple}\fi
 \node[v3 title] at (139,\y+8) {\out};
 \draw[v3 flow] (42.4,\y-3)--(83.62,\y-3);
 \ifnum\y=43
  \draw[v3 flow] (92.38,\y-3)--(135.01,\y-3);
 \else
  \draw[v3 flow] (92.38,\y-3)--(131.75,\y-3);
 \fi
}
\fill[FigurePurple,rounded corners=.4pt] (39,13) rectangle ++(2.5,2.5);
\draw[v3 flow,dashed] (125,77.5)--(125,59)--(88,59)--(88,46.1);
\draw[v3 flow,dashed] (88,33.75)--(88,23.1);
\node[v3 note,anchor=west] at (90,61.5) {共享参数 $\theta$};
\node[v3 note] at (81,1) {干净输入上的性能不能排除后门；此处仅示意条件性异常};
\end{tikzpicture}
